\section{What an ordinal statement means in iPRA}
\label{sec:foundation}
The article \emph{Ordinal number} treats ordinals both as transitive sets and as order types. Two examples suggest how to read it arithmetically. The assertion $3=\{0,1,2\}$ can be a literal calculation with four integers coding finite sets. The assertion $\omega^2=\sup_n\omega n$ has a different reading. Every point of $\omega^2$ lies in a finite-numbered block of type $\omega$. Given a point, we can compute the number of a later block. We work out both readings, including the maps and witnesses which distinguish them. The final section gives the source correspondence.

\subsection{Induction and finite data}
Our iPRA is first-order intuitionistic logic with equality, the zero and successor axioms, a symbol and defining equations for each primitive recursive (PR) function, and induction for quantifier-free formulas with numerical parameters. A PR definition is built by composition and ordinary numerical primitive recursion. We use $\mathbb N$ to include zero. Searches below a supplied bound return a first index, or a default outside the range. A predicate whose characteristic function is PR is decidable here. Trichotomy for a coded comparison permits a case split even when the conclusion of the argument contains quantifiers.

Write $\mathrm{All}(0,p)=1$ and $\mathrm{All}(s+1,p)=\mathrm{All}(s,p)\cdot p(s)$ for a fixed Boolean PR test $p$ with arbitrary fixed numerical parameters. A least-failure register in the same loop proves that the output is $1$ exactly when every index below the bound passes. If the output is $0$, the register supplies a failing index. The analogous disjunction scan supplies a first witness. These claims use the loop equations and induction on numerical equalities or bounded tests. For example, the saved failure flag implies $j<s$ and $p(j)=0$, while the forward direction for a selected $i<s$ follows by carrying that fixed index through the loop. Nested bounded quantifiers are consequently shorthand for a single PR test. An induction below using a bounded universal uses that test as induction formula. We introduce universal input variables after proving the instance with them as parameters. We use ordinary intuitionistic quantifiers for conclusions, for hypotheses available at particular arguments, and for schemata over fixed PR definitions. There is no universal PR evaluator for all PR programs.

Here is one consistent realization of the data structures. Code an empty list by $0$ and a cons cell with head $u$ and tail $v$ by
\[
 C(u,v)=1+(u+v)^2+u .
\]
The two components are smaller than this positive code. To decode $c>0$, set $s=\lfloor\sqrt{c-1}\rfloor$ by bounded search, $u=(c-1)-s^2$. Require $u\le s$, and set $v=s-u$. In the permitted range $s^2+u<(s+1)^2$, so this is inverse to construction. Test the pair again on decoding, and use fixed defaults on malformed data. Valid binary trees have valid subtrees in both positions. A list with arbitrary numerical entries needs only its tail to be a valid list. Scans, length, retrieval, concatenation and insertion sort follow by recursion with at most the input code (or the explicit length) many steps. To tabulate $f(t)$ with calls on smaller numerical codes, retain the list $[f(0),\ldots,f(t-1)]$ and retrieve each requested entry by a bounded scan. The history itself is built by one PR recursion. Its potentially large numerical code is allowed. This also implements a fixed finite product of recursions by packing the registers. The same construction works for pairs and triples of subtrees by loops over their decreasing total code or by tables over a fixed box of input codes.

\subsection{The notation and its order}
A term is a finite tree displayed as the list of its children. Compare two lists lexicographically using this same comparison on entries, with the proper prefix smaller. Let $\mathcal O$ consist of the trees for which every list, including every descendant list, is weakly decreasing. Empty is written $0$. Write $[a]$ for a singleton, $1=[0]$, $\omega=[1]$. An expression such as $[2,2,0]$ denotes $\omega^2\cdot 2+1$. Repetitions carry coefficients. The exponent of the largest power in a nonempty term is its first child, also called its degree. A code for a sequence of terms and the list of children of a single term can have the same numerical representation. The indicated domain test specifies which is meant.

\begin{theorem}[Comparison, normal forms, and finite extrema]\label{thm:order}
Validity, equality and comparison are PR. Comparison gives a strict linear order with least element $0$. Each term is uniquely specified by the strictly descending list of its distinct exponents and their positive multiplicities. Every nonempty coded finite family of terms has a witnessed least and greatest member. Its least upper bound is its greatest member. The empty family has least upper bound $0$.
\end{theorem}
\begin{proof}
Compare the empty and nonempty cases directly. For two nonempty cells compare the heads and, only if equal, the tails. The recursive pairs have smaller total code and remain inside the input box. Thus a PR table computes the comparison, including the equality output. A second table tests validity using the two subtrees and the order of adjacent heads. Equality output forces equality of both heads and tails by the same decreasing computation, whence equality of codes. The converse is immediate by the defining clauses. Interchanging inputs reverses the output. For transitivity of $x<y<z$ write their respective heads as $a,b,c$ when all are nonempty. The two comparisons give $a\le b\le c$. If either is strict, transitivity on smaller heads gives $a<c$ (or substitution does so in an equality case). If both are equal, the two strict tail comparisons give the result on smaller tails. The empty cases use the prefix rule. Formally fix a numerical bound $B$ on the three original inputs. Induct on $s$ using the Boolean conjunction of the comparison implications for all triples with components $\le B$ and total code $\le s$. Every subcomparison used in the proof is covered at the preceding step. This is quantifier-free induction, independent of the putative well-ordering of the terms.

Grouping equal adjacent exponents is inverted by repeating each exponent its recorded number of times. A first differing exponent or multiplicity decides comparison, so distinct grouped normal forms remain distinct. A scan of a finite family updates the index of its current extremum. The invariant says its stored value is an input and bounds all scanned entries. For a least upper bound, every upper bound bounds the selected greatest entry. For the empty list every term is an upper bound and $0$ is least.
\end{proof}

\paragraph{Internal and external readings.}
Internally $x<a$ is a decidable numerical predicate on valid codes. We write $I_a$ for this predicate, the strict initial segment of $a$. An isomorphism of displayed presentations means explicit mutually inverse maps preserving the displayed comparisons. A supremum statement may be a universal property quantified over codes. Its proof below usually computes a witness against every smaller candidate. Externally the conventional interpretation of a tree is $|[a_1,\ldots,a_k]|=\omega^{|a_1|}+\cdots+\omega^{|a_k|}$. Classical Cantor normal form identifies precisely the ordinals below $\varepsilon_0$ with these trees. To recall the range argument, the finite heights of trees bound their values by finite towers of $\omega$, and every ordinal below such a tower has a finite CNF whose exponents are below the previous tower. Induction on the tower height represents it. Classical uniqueness and the usual lexicographic comparison give injectivity. This paragraph identifies the classical referents. The arithmetical proofs below work on codes directly. All unqualified term variables from here on satisfy the validity test.
